% ===========================================================================
%  樊启斌《高等代数典型问题与方法》 · 第 9 章 欧氏空间
%  §9.2 正交矩阵与正交变换（印刷页 595–607）+ §9.4 Gram 矩阵（印刷页 628–634）
%  —— 深化提取（服务于考研大纲〔向量〕：内积、规范正交基、正交矩阵）
%
%  源：PDF p604–p616 = 印刷页 595–607（§9.2，实测页眉页码吻合 +9）
%      PDF p637–p643 = 印刷页 628–634（§9.4 Gram 矩阵）
%      偏移 PDF = 印刷页 + 9，与 AGENTS.md §3.4 一致。
%
%  分工：本片段只做「正交矩阵的判定与性质」「正交变换保内积/保长度」
%        「Gram 矩阵及其正定性」这三块偏计算与判定面向的内容；
%        §5.1 线性空间与子空间、§9.1 内积与欧氏空间的定义性内容由另一分片负责。
%
%  ⚠ 与主控手写块 deep_ch06_authored.tex 的分工：那个块已给出正交矩阵的六条
%    等价刻画（$Q^{\mathrm T}Q=E$、$QQ^{\mathrm T}=E$、列/行为规范正交组、
%    $Q^{-1}=Q^{\mathrm T}$、保长度）。本片段**不重复**这六条，只补它没有的
%    性质与判定细节：特征值模为 1、$|Q|=\pm1$ 的分类、代数余子式、上三角情形、
%    从等式判定正交、Gram 矩阵及其正定性。
%
%  提取原则：只取定义 / 定理 / 性质 / 重要结论 / 方法要点 / 关键证明思路；
%            例题的数值演算过程不进正文，只在承载了可复用结论时提取
%            「结论 + 一句话证明思路」。正交变换的完整分类理论（标准形）、
%            QR 分解 / Givens / Householder 旋转按考纲外跳过（明细见文件末）。
% ===========================================================================


% ---------------------------------------------------------------------------
% 【深化】服务考点：正交变换的等价刻画（考纲〔向量〕正交矩阵的概念与性质）
% 来源：樊 §9.2「基本理论与要点提示」9.2.2，印刷页 595（PDF p604）
% ---------------------------------------------------------------------------
\begin{fdeep}{正交变换：四条等价刻画（把矩阵语言升级为变换语言）}
欧氏空间 $V$ 的\textbf{线性变换} $\mathscr{A}$，下面四条命题相互等价：
（1）$\mathscr{A}$ 是 $V$ 的正交变换，即\textbf{保持内积}：$\forall\alpha,\beta\in V$，$(\mathscr{A}(\alpha),\mathscr{A}(\beta))=(\alpha,\beta)$；
（2）$\mathscr{A}$ 在 $V$ 的任一标准正交基下的矩阵为正交矩阵；
（3）$\mathscr{A}$ 将 $V$ 的标准正交基变为标准正交基；
（4）$\mathscr{A}$ 保持 $V$ 中所有向量的长度：$\forall\alpha\in V$，$|\mathscr{A}(\alpha)|=|\alpha|$。
\textbf{它比矩阵刻画多说了什么}：$Q^{\mathrm T}Q=E$ 只说明"矩阵的列是规范正交组"，（1）才说明"变换不改变任意两个向量的内积"——这才是正交变换的本质，（4）只是（1）取 $\alpha=\beta$ 的特例。另注意（2）中是"任一"：只要在某一个标准正交基下矩阵是正交阵，它在所有标准正交基下就都 是正交阵（两组标准正交基之间的过渡矩阵本身是正交阵）。

【反哺点】服务考纲〔向量〕"了解规范正交基、正交矩阵的概念以及它们的性质"。选择题问"下列哪些条件等价于 $\mathscr{A}$ 为正交变换"时这四条通常全选；真正的陷阱在"保长度"三个字：矩阵层面 $|Qx|=|x|$ 与 $Q^{\mathrm T}Q=E$ 等价，变换层面却必须配上\textbf{线性}才等价（见下一块）。
\end{fdeep}


% ---------------------------------------------------------------------------
% 【深化】服务考点：正交变换的判定（考纲〔向量〕正交矩阵）
% 来源：樊 例 9.47，印刷页 606（PDF p615）
% ---------------------------------------------------------------------------
\begin{fdeep}{保持内积的变换自动是线性的}
设 $\sigma$ 是 $n$ 维欧氏空间 $V$ 的\textbf{变换}（不预设线性），若 $\sigma$ 保持内积不变，即 $\forall\alpha,\beta\in V$ 有 $(\sigma\alpha,\sigma\beta)=(\alpha,\beta)$，则 $\sigma$ 一定是线性变换，从而是正交变换。
\textbf{证明思路}（两次"长度为零"的展开）：由保内积，
\fml{\bigl(\sigma(\alpha+\beta)-\sigma\alpha-\sigma\beta,\ \sigma(\alpha+\beta)-\sigma\alpha-\sigma\beta\bigr)
=(\alpha+\beta,\alpha+\beta)-2(\alpha+\beta,\alpha)-2(\alpha+\beta,\beta)+(\alpha,\alpha)+(\beta,\beta)+2(\alpha,\beta)=0,}
故 $\sigma(\alpha+\beta)=\sigma\alpha+\sigma\beta$；同理
$\bigl(\sigma(k\alpha)-k\sigma\alpha,\sigma(k\alpha)-k\sigma\alpha\bigr)=(k\alpha,k\alpha)-2k(k\alpha,\alpha)+k^{2}(\alpha,\alpha)=0$，故 $\sigma(k\alpha)=k\sigma\alpha$。

【反哺点】服务考纲〔向量〕"了解内积……正交矩阵"。若题目只给"$\sigma$ 保持内积"，不必怀疑"没给线性所以条件不足"——线性是免费得到的；这条也解释了为什么"保内积"可以单独作为正交变换的定义。
\end{fdeep}


% ---------------------------------------------------------------------------
% 【深化】服务考点：正交变换的判定（考纲〔向量〕正交矩阵）
% 来源：樊 例 9.48，印刷页 606（PDF p615）
% ---------------------------------------------------------------------------
\begin{fdeep}{保距变换还差一个"保原点"}
设 $\sigma$ 是 $n$ 维欧氏空间 $V$ 的保距变换，即 $|\sigma(\alpha)-\sigma(\beta)|=|\alpha-\beta|\ (\forall\alpha,\beta\in V)$，则\textbf{只有}当 $\sigma(0)=0$ 时才能推出 $\sigma$ 是正交变换。
\textbf{为什么}：保距只保证"任意两点的距离不变"，只有 $\sigma(0)=0$ 时才有
\fml{|\sigma(\alpha)|=|\sigma(\alpha)-\sigma(0)|=|\alpha-0|=|\alpha|,}
即"保距"过渡到"保长度"必须靠 $\sigma(0)=0$ 这座桥；再把保距条件按内积展开
$(\sigma\alpha-\sigma\beta,\sigma\alpha-\sigma\beta)=(\alpha-\beta,\alpha-\beta)$，即得 $(\sigma\alpha,\sigma\beta)=(\alpha,\beta)$。
\textbf{反例}：取定 $\xi\ne0$，令 $\sigma(\alpha)=\alpha+\xi$。它显然不是线性变换，却满足 $|\sigma(\alpha)-\sigma(\beta)|=|\alpha-\beta|$，所以 $\sigma(0)=0$ 不能省略。

【反哺点】服务考纲〔向量〕"了解内积"。这是"保长度"与"保距离"最容易混的一处：保长度（$|\sigma\alpha|=|\alpha|$）本身等价于正交，保距离则多出一个平移自由度。考场上见到"$|\sigma(\alpha)-\sigma(\beta)|=|\alpha-\beta|$"务必先验 $\sigma(0)=0$。
\end{fdeep}


% ---------------------------------------------------------------------------
% 【深化】服务考点：正交矩阵的特征值（考纲〔特征值与特征向量〕）
% 来源：樊 例 9.34，印刷页 597（PDF p606）
% ---------------------------------------------------------------------------
\begin{fdeep}{正交矩阵的特征值模恒为 $1$}
设 $\lambda$ 是 $n$ 阶正交矩阵 $A$ 的特征值，$\alpha\ne0$ 是相应的特征向量，则 $A\alpha=\lambda\alpha$ 蕴含 $|\lambda|=1$。证明只需一步配对：对共轭后的等式 $A\bar\alpha=\bar\lambda\bar\alpha$ 取转置，再与自身配对，
\fml{\bar\alpha^{\mathrm T}\alpha=\bar\alpha^{\mathrm T}A^{\mathrm T}A\alpha=(A\bar\alpha)^{\mathrm T}(A\alpha)=|\lambda|^{2}\,\bar\alpha^{\mathrm T}\alpha,}
而 $\bar\alpha^{\mathrm T}\alpha\ne0$，故 $|\lambda|^{2}=1$。
\textbf{一条常考的加强结论}：若 $\lambda$ 为虚数，把特征向量写成 $\alpha=\beta+\mathrm i\gamma$（$\beta,\gamma\in\mathbb R^{n}$，$\gamma\ne0$）。由上式知 $\alpha$ 与 $\bar\alpha$ 是 $A$ 的属于\textbf{不同}特征值 $\lambda,\bar\lambda$ 的特征向量，而正交矩阵属于不同特征值的特征向量必正交，故 $(\alpha,\bar\alpha)=0$，即 $(\beta,\beta)-(\gamma,\gamma)+\mathrm i\bigl[(\beta,\gamma)+(\gamma,\beta)\bigr]=0$。分开实部与虚部得 $(\beta,\gamma)=0$ 且 $|\beta|=|\gamma|$：\textbf{复特征向量的实部与虚部必正交且等长}。

【反哺点】服务考纲〔特征值〕"掌握特征值和特征向量的概念与求法"及〔向量〕正交矩阵。看到"正交矩阵"就立刻写两条：特征值模为 $1$；属于不同特征值的特征向量正交。前者可秒判"某数能否为正交阵的特征值"。
\end{fdeep}


% ---------------------------------------------------------------------------
% 【深化】服务考点：$|Q|=\pm1$ 推出的结论（考纲〔行列式〕〔特征值〕〔向量〕）
% 来源：樊 例 9.29，印刷页 595（PDF p604）；例 9.35(1)，印刷页 597
% ---------------------------------------------------------------------------
\begin{fdeep}{$|Q|=\pm1$ 的两条直接推论}
\textbf{推论一}：设 $A$ 是 $n$ 阶正交矩阵且 $|A|=-1$，则 $-1$ 必是 $A$ 的特征值。证明只用行列式的乘法与转置：
\fml{|E+A|=|AA^{\mathrm T}+A|=|A|\,|A^{\mathrm T}+E|=|A|\,|A+E|=-|E+A|,}
故 $|E+A|=0$。把它与"奇数阶 $|A|=1$"合起来可统一记忆：正交矩阵的实特征值只能是 $\pm1$；$n$ 为奇数时复特征值成共轭对出现、乘积为 $1$，故实特征值有奇数个且乘积 $=|A|$，于是 $|A|=1$ 时 $1$ 必是特征值（$|A|=-1$ 时 $-1$ 必是特征值）。三阶正交阵中这是最常用的一条。
\textbf{推论二}：设 $A,B$ 都是正交矩阵且 $|A|+|B|=0$，则 $|A+B|=0$。证明：由 $|A|^{2}=|B|^{2}=1$ 与 $|A|+|B|=0$ 得 $|AB|=-1$，于是
\fml{|A+B|=|AB^{\mathrm T}B+BA^{\mathrm T}B|=|A|\,|(A+B)^{\mathrm T}|\,|B|=-|A+B|.}

【反哺点】服务考纲〔向量〕正交矩阵、〔行列式〕行列式的性质。"正交 + 行列式为 $-1$（或奇数阶且行列式为 $1$）$\Rightarrow$ 必有 $\pm1$ 特征值"是选择题与证明题的高频切口；而 $|A|+|B|=0$ 推出 $|A+B|=0$ 的套路就是"乘 $A^{\mathrm T}A$ 或 $B^{\mathrm T}B$ 造出转置"。
\end{fdeep}


% ---------------------------------------------------------------------------
% 【深化】服务考点：3 阶正交矩阵的特征多项式（考纲〔特征值〕〔二次型〕）
% 来源：樊 例 9.36，印刷页 598（PDF p607）
% ---------------------------------------------------------------------------
\begin{fdeep}{三阶正交矩阵的特征多项式必为 $\lambda^{3}-a\lambda^{2}+a\lambda-1$}
若 $A$ 是三阶正交矩阵且 $|A|=1$，则
\fml{f(\lambda)=|\lambda E-A|=\lambda^{3}-a\lambda^{2}+a\lambda-1,}
其中 $a$ 是某个实数。推导：设 $f(\lambda)=\lambda^{3}+c_{1}\lambda^{2}+c_{2}\lambda+c_{3}$，由 $c_{3}=f(0)=|{-}A|=(-1)^{3}|A|=-1$，再用"$1$ 是特征值"得 $f(1)=1+c_{1}+c_{2}+c_{3}=0$，即 $c_{1}+c_{2}=0$，令 $c_{2}=a$、$c_{1}=-a$ 即得。
\textbf{系数 $a$ 的取值范围}：另两个特征值 $\lambda_{2},\lambda_{3}$ 满足 $|\lambda_{2}|=|\lambda_{3}|=1$，由根与系数关系
\fml{a=\lambda_{1}+\lambda_{2}+\lambda_{3}=1+\lambda_{2}+\lambda_{3},}
再由 $|\lambda_{2}+\lambda_{3}|\le|\lambda_{2}|+|\lambda_{3}|=2$ 得 $-1\le a\le3$。

【反哺点】服务考纲〔特征值〕〔二次型〕。三阶正交阵的题目几乎都能用这一步解决：不必展开行列式，直接设特征多项式为 $\lambda^{3}-a\lambda^{2}+a\lambda-1$，其中 $a=\operatorname{tr}A$，再用 $\pm1\le a\le3$ 定参数范围——这也解释了为什么考题总把待定参数设在迹上。
\end{fdeep}


% ---------------------------------------------------------------------------
% 【深化】服务考点：3 阶正交矩阵的迹（考纲〔特征值〕〔行列式〕）
% 来源：樊 例 9.40，印刷页 600（PDF p609）；例 9.35【注 1】，印刷页 598
% ---------------------------------------------------------------------------
\begin{fdeep}{三阶正交矩阵的迹：一个恒等式与旋转角公式}
\textbf{恒等式}：若 $A$ 是三阶正交矩阵且 $|A|=-1$，则 $\operatorname{tr}A^{2}=2\operatorname{tr}A+(\operatorname{tr}A)^{2}$。
验证（也是通用证法）：$|A|=-1$ 时三个特征值为 $-1,\mathrm e^{\mathrm i\varphi},\mathrm e^{-\mathrm i\varphi}$，于是 $\operatorname{tr}A=-1+2\cos\varphi$，$\operatorname{tr}A^{2}=1+2\cos2\varphi=4\cos^{2}\varphi-1$，而 $2\operatorname{tr}A+(\operatorname{tr}A)^{2}=2(2\cos\varphi-1)+(2\cos\varphi-1)^{2}=4\cos^{2}\varphi-1$，两边相等。
\textbf{旋转角公式}：三阶正交阵 $A=(a_{ij})$ 且 $\det A=1$ 时
\fml{\varphi=\arccos\frac{a_{11}+a_{22}+a_{33}-1}{2},}
即迹直接给出旋转角 $\varphi$（几何意义：绕过原点的直线的旋转）。

【反哺点】服务考纲〔特征值〕〔行列式〕。这两条是"迹与特征值"关系的标准演练：只要题目给出三阶正交阵，就可以把 $\operatorname{tr}A$、$\operatorname{tr}A^{2}$、$|A|$ 三者用 $\varphi$ 一个参数串起来，含 $\operatorname{tr}$ 的恒等式题与含 $\varphi$ 的求值题都能一步到位，不用碰矩阵元素。
\end{fdeep}


% ---------------------------------------------------------------------------
% 【深化】服务考点：正交矩阵的元素与其代数余子式（考纲〔行列式〕〔矩阵〕）
% 来源：樊 例 9.31，印刷页 596（PDF p605）
% ---------------------------------------------------------------------------
\begin{fdeep}{正交矩阵的元素与代数余子式：$a_{ij}=\pm A_{ij}$}
设 $A=(a_{ij})$ 是 $n$ 阶\textbf{实}方阵，下面两组命题等价：
（一）$A$ 是正交矩阵且 $|A|=1$ 或 $|A|=-1$；
（二）$|A|=1$ 时，$A$ 的每一个元素等于该元素的代数余子式，即 $a_{ij}=A_{ij}$；$|A|=-1$ 时，$A$ 的每一个元素等于该元素的代数余子式的相反数，即 $a_{ij}=-A_{ij}$（$1\le i,j\le n$）。
\textbf{证明思路（一句话）}：核心是伴随矩阵的恒等式 $AA^{*}=|A|E$ 与 $A^{*}=|A|A^{-1}$；正交时 $A^{-1}=A^{\mathrm T}$，故 $A^{*}=|A|A^{-1}=|A|A^{\mathrm T}$，把 $|A|=1$ 或 $-1$ 代入并逐元素比较即得。反过来只用到 $A^{*}=\pm A^{\mathrm T}$ 与 $AA^{*}=|A|E$，直接算出 $A^{\mathrm T}A=E$。

【反哺点】服务考纲〔行列式〕代数余子式、〔矩阵〕伴随矩阵与逆矩阵。这是一个\textbf{反着用}的判定：题目若给出"每个元素等于它的代数余子式"（或都等于代数余子式的相反数），不必展开九个行列式，直接断言 $A$ 正交且 $|A|=\pm1$；反之已知正交，即可把元素与代数余子式互相替换，用来化简含 $A^{*}$ 的式子。
\end{fdeep}


% ---------------------------------------------------------------------------
% 【深化】服务考点：正交矩阵的结构判定（考纲〔矩阵〕〔向量〕）
% 来源：樊 例 9.41，印刷页 600–601（PDF p609–p610）
% ---------------------------------------------------------------------------
\begin{fdeep}{上（下）三角的正交矩阵必为对角矩阵，对角元只能是 $\pm1$}
\textbf{结论}：任一上三角正交矩阵必为对角矩阵，且对角线上的元素为 $1$ 或 $-1$；下三角情形同理。
\textbf{证明思路一（用求逆的封闭性）}：$A$ 上三角正交，则 $A^{-1}$ 也是上三角；而 $A^{-1}=A^{\mathrm T}$ 是下三角，既上三角又下三角的矩阵必为对角矩阵；再由 $A^{\mathrm T}A=E$ 得 $a_{ii}^{2}=1$，故 $a_{ii}=\pm1$。
\textbf{证明思路二（分块归纳）}：$n=1$ 显然。设结论对 $n-1$ 阶成立，把 $A$ 分块为 $A=\begin{pmatrix}a_{11}&\alpha\\ 0&B\end{pmatrix}$，其中 $\alpha$ 为 $1\times(n-1)$ 行向量、$B$ 为 $n-1$ 阶上三角矩阵。由
\fml{A^{\mathrm T}A=\begin{pmatrix}a_{11}&0\\ \alpha^{\mathrm T}&B^{\mathrm T}\end{pmatrix}
\begin{pmatrix}a_{11}&\alpha\\ 0&B\end{pmatrix}
=\begin{pmatrix}a_{11}^{2}&a_{11}\alpha\\ a_{11}\alpha^{\mathrm T}&\alpha^{\mathrm T}\alpha+B^{\mathrm T}B\end{pmatrix}=E}
得 $a_{11}^{2}=1$（即 $a_{11}=\pm1$）、$a_{11}\alpha=0$（故 $\alpha=0$）、$B^{\mathrm T}B=E_{n-1}$。可见 $B$ 是 $n-1$ 阶上三角正交矩阵，由归纳假设 $B$ 为对角阵且对角元为 $\pm1$，从而 $A$ 也是如此。

【反哺点】服务考纲〔矩阵〕〔向量〕正交矩阵的判定。选择题里几乎秒杀："可能是正交矩阵"的选项中，非对角的上三角矩阵一律排除。它也是 $A=QR$ 分解唯一性的关键引理——两个上三角正交因子相比只能得到对角矩阵。
\end{fdeep}


% ---------------------------------------------------------------------------
% 【深化】服务考点：正交矩阵与对称矩阵的关系（考纲〔相似理论〕〔二次型〕）
% 来源：樊 例 9.44，印刷页 603（PDF p612）
% ---------------------------------------------------------------------------
\begin{fdeep}{正交矩阵的特征值全为实数 $\iff$ 它是对称矩阵}
$n$ 阶正交矩阵 $A$ 的特征值全为实数，当且仅当 $A$ 是对称矩阵；此时 $A$ 的特征值只能取 $\pm1$。
\textbf{证明思路}：$n$ 阶实方阵 $A$ 为正交矩阵当且仅当存在正交矩阵 $Q$ 使 $Q^{\mathrm T}AQ=\operatorname{diag}(E_{r},-E_{s},A_{1},A_{2},\dots,A_{t})$，其中 $A_{j}$ 是二阶旋转块，而
\fml{|\lambda E-A_{j}|=(\lambda-\cos\theta_{j})^{2}+\sin^{2}\theta_{j}}
只有当 $\sin\theta_{j}=0$ 时才有实根。题设"特征值全为实数"迫使 $\sin\theta_{j}=0$，此时每个 $A_{j}$ 都是 $\pm E$ 型对角块，于是 $Q^{\mathrm T}AQ=D$ 是对角矩阵（对角元只取 $\pm1$），故 $A=QDQ^{\mathrm T}$ 对称；反之对称矩阵特征值必为实数。
若只做判定，可记更强的等价链：$A^{\mathrm T}=A$ 且 $Q^{\mathrm T}Q=E$ $\iff$ 特征值全为 $\pm1$ $\iff$ 存在正交 $Q$ 使 $A=Q\operatorname{diag}(\pm1)Q^{\mathrm T}$。

【反哺点】服务考纲〔相似理论〕"实对称矩阵必可正交相似对角化"与〔二次型〕。它把"正交"这个不变量与"对称"这个形状接了起来：考场上遇到"正交 + 对称"的条件组合，直接写特征值为 $\pm1$、$A^{2}=E$，可以省掉大量推导。
\end{fdeep}


% ---------------------------------------------------------------------------
% 【深化】服务考点：从等式判定正交（考纲〔向量〕正交矩阵）
% 来源：樊 例 9.38，印刷页 599（PDF p608）
% ---------------------------------------------------------------------------
\begin{fdeep}{判定补充一：$(A+I)^{-1}+(A^{\mathrm T}+I)^{-1}=I$}
设 $A$ 是 $n$ 阶实方阵，且特征值不等于 $-1$，则 $A+I$ 与 $A^{\mathrm T}+I$ 都可逆，并且
\fml{A\ \text{是正交矩阵}\iff (A+I)^{-1}+(A^{\mathrm T}+I)^{-1}=I.}
\textbf{证明要点}："$A+I$ 可逆"来自"$-1$ 不是特征值"，即 $|{-}I-A|=(-1)^{n}|I+A|\ne0$（同理 $|A^{\mathrm T}+I|=|A+I|\ne0$）。充分性从等式出发，两边左乘 $A+I$ 得 $I+(A+I)(A^{\mathrm T}+I)^{-1}=A+I$，再右乘 $A^{\mathrm T}+I$，注意 $(A^{\mathrm T}+I)^{\mathrm T}=A+I$，即得 $AA^{\mathrm T}=I$；必要性把上述过程倒着走一遍。

【反哺点】服务考纲〔向量〕"了解正交矩阵"与〔矩阵〕逆矩阵。考纲内的正交判定一旦不是直接给 $Q^{\mathrm T}Q=E$，最常见的一类就是这种\textbf{含 $(A+I)^{-1}$ 的求逆等式}：见到它不要去算逆，直接按"左乘 $A+I$、右乘 $A^{\mathrm T}+I$"两步消项，两三行就出 $AA^{\mathrm T}=I$。
\end{fdeep}


% ---------------------------------------------------------------------------
% 【深化】服务考点：从等式判定正交（考纲〔矩阵〕迹、〔向量〕正交矩阵）
% 来源：樊 例 9.45，印刷页 603（PDF p612）
% ---------------------------------------------------------------------------
\begin{fdeep}{判定补充二：元素平方和 $\sigma(A)=\operatorname{tr}(A^{\mathrm T}A)$ 的恒等式}
记 $A=(a_{ij})$ 的全部元素的平方和为 $\sigma(A)=\sum_{i=1}^{n}\sum_{j=1}^{n}a_{ij}^{2}$，则
（1）$\sigma(A)=\operatorname{tr}(A^{\mathrm T}A)$；
（2）$\sigma(A^{\mathrm T}A)=\sigma(AA^{\mathrm T})$；
（3）$A$ 为正交矩阵的充分必要条件是：对任意 $n$ 阶实方阵 $B$，都有 $\sigma(A^{\mathrm T}BA)=\sigma(B)$。
\textbf{（3）的证明要点}：必要性直接算
$\sigma(A^{\mathrm T}BA)=\operatorname{tr}\bigl((A^{\mathrm T}BA)^{\mathrm T}(A^{\mathrm T}BA)\bigr)=\operatorname{tr}(B^{\mathrm T}B)=\sigma(B)$；充分性取 $B=E_{ii}$ 得 $\sigma(A^{\mathrm T}E_{ii}A)=1$，逐列算出各列模长为 $1$，再取 $B=E$ 得 $\sigma(A^{\mathrm T}A)=n$，两者结合推出 $A^{\mathrm T}A=E$。

【反哺点】服务考纲〔向量〕"了解正交矩阵"与〔矩阵〕迹的计算。$\sigma(A)=\operatorname{tr}(A^{\mathrm T}A)=\sum a_{ij}^{2}$ 这条把"元素平方和、迹、内积（列向量模长平方之和）"三者打通，是含 $\operatorname{tr}(A^{\mathrm T}A)$ 类题目的统一入口；正定判定中常用的"对角元之和 = 全部元素平方和"也是它。
\end{fdeep}


% ---------------------------------------------------------------------------
% 【深化】服务考点：正交变换的求法（考纲〔向量〕正交变换/规范正交基）
% 来源：樊 例 9.49，印刷页 606–607（PDF p615–p616）；例 9.50，印刷页 607
% ---------------------------------------------------------------------------
\begin{fdeep}{求正交变换的三条现成路径}
\textbf{路径一（最通用：指定基的像，补全正交矩阵）}。设 $\eta_{1},\eta_{2},\eta_{3}$ 是规范正交基，令 $\sigma(\eta_{i})=x_{i1}\eta_{1}+x_{i2}\eta_{2}+x_{i3}\eta_{3}$，则 $\sigma$ 在该基下的矩阵 $A$ 以像的坐标为列。要求 $\sigma$ 为正交变换，只需 $A$ 是正交矩阵，于是\textbf{只须凑出最后一列}：解"第三列与已知两列正交"的两个线性方程，再单位化即可。两个解（差一个符号）都合法。
\textbf{路径二（最直接：设矩阵解方程组）}。设 $\sigma$ 在标准正交基下的矩阵为 $A$，由 $A^{\mathrm T}A=E$ 建立方程组，与已知的像条件联立求解。
\textbf{路径三（最几何：两组标准正交基互换）}。把已知向量组用施密特方法分别补成两组标准正交基 $\xi_{1},\dots,\xi_{n}$ 与 $\eta_{1},\dots,\eta_{n}$，再定义线性变换 $\sigma(\xi_{i})=\eta_{i}\ (i=1,\dots,n)$；由"标准正交基的像仍是标准正交基"知 $\sigma$ 是正交变换且已满足题设。

【反哺点】服务考纲〔向量〕"了解规范正交基""了解正交矩阵……的性质"。考研中"求一个正交变换"的题几乎都落在路径一或路径二：路径一的好处是\textbf{只需凑出最后一列}（用正交条件解两个方程），比整块验证 $A^{\mathrm T}A=E$ 快得多；路径三则在"已知两组向量"的存在性题里最好用。
\end{fdeep}


% ---------------------------------------------------------------------------
% 【深化】服务考点：Gram 矩阵的定义与来源（考纲〔向量〕内积）
% 来源：樊 §9.4.1，印刷页 628（PDF p637）
% ---------------------------------------------------------------------------
\begin{fdeep}{Gram 矩阵（度量矩阵）：定义与三种来源}
\textbf{定义}：欧氏空间 $V$ 中向量组 $\alpha_{1},\alpha_{2},\dots,\alpha_{m}$ 的 Gram 矩阵为
\fml{G(\alpha_{1},\alpha_{2},\dots,\alpha_{m})=
\begin{pmatrix}
(\alpha_{1},\alpha_{1})&(\alpha_{1},\alpha_{2})&\cdots&(\alpha_{1},\alpha_{m})\\
(\alpha_{2},\alpha_{1})&(\alpha_{2},\alpha_{2})&\cdots&(\alpha_{2},\alpha_{m})\\
\vdots&\vdots&&\vdots\\
(\alpha_{m},\alpha_{1})&(\alpha_{m},\alpha_{2})&\cdots&(\alpha_{m},\alpha_{m})
\end{pmatrix}.}
欧氏空间 $V$ 中关于基 $\alpha_{1},\dots,\alpha_{n}$ 的度量矩阵就是 $G(\alpha_{1},\dots,\alpha_{n})$；特别地，$V$ 的任一标准正交基的度量矩阵是单位矩阵 $E$。
\textbf{第三种来源（最实用）}：若向量组 $\alpha_{1},\dots,\alpha_{m}$ 在标准正交基下的坐标分别为矩阵 $A$ 的列向量，则
\fml{G(\alpha_{1},\alpha_{2},\dots,\alpha_{m})=A^{\mathrm T}A.}

【反哺点】服务考纲〔向量〕"了解内积的概念"。这是把"内积条件"翻译成\textbf{矩阵方程}的桥梁：凡题目出现"$(\beta,\alpha_{i})=b_{i}$"，写出的系数矩阵就是 Gram 矩阵；另外 $A^{\mathrm T}A$、$AA^{\mathrm T}$ 一律要认成 Gram 矩阵（分别是列、行向量组的），这是把抽象内积题"矩阵化"的第一步。
\end{fdeep}


% ---------------------------------------------------------------------------
% 【深化】服务考点：Gram 矩阵的秩（考纲〔向量〕内积、线性相关性）
% 来源：樊 例 9.88、例 9.89，印刷页 628–629（PDF p637–p638）
% ---------------------------------------------------------------------------
\begin{fdeep}{Gram 矩阵的秩 $=$ 向量组的秩（线性无关判据）}
\fml{\operatorname{rank}G(\alpha_{1},\dots,\alpha_{m})=\operatorname{rank}(\alpha_{1},\dots,\alpha_{m});}
特别地，$\alpha_{1},\dots,\alpha_{m}$ 线性无关当且仅当它们的 Gram 矩阵非奇异（$m$ 阶方阵时即 $|G|\ne0$）。
\textbf{证明要点}（把抽象内积变成矩阵运算）：把每个 $\alpha_{j}$ 用标准正交基 $\eta_{1},\dots,\eta_{n}$ 展开得 $\alpha_{j}=\sum_{k=1}^{n}b_{kj}\eta_{k}$，代入内积得
\fml{(\alpha_{i},\alpha_{j})=\Bigl(\sum_{k=1}^{n}b_{ki}\eta_{k},\sum_{k=1}^{n}b_{kj}\eta_{k}\Bigr)=\sum_{k=1}^{n}b_{ki}b_{kj},}
即 $G=B^{\mathrm T}B$（$B=(b_{ij})_{n\times m}$）。于是 $\operatorname{rank}G=\operatorname{rank}(B^{\mathrm T}B)=\operatorname{rank}B=\operatorname{rank}(\alpha_{1},\dots,\alpha_{m})$。

【反哺点】服务考纲〔向量〕"了解内积"与〔向量〕线性相关性的判定。它给出一条与坐标系无关的线性无关判据：不必解方程组，只算 Gram 行列式是否为零；反过来，抽象空间里"内积为定值"的线性方程组，其系数矩阵就是 Gram 矩阵，可逆即唯一解。
\end{fdeep}


% ---------------------------------------------------------------------------
% 【深化】服务考点：Gram 矩阵的正定性（考纲〔二次型〕正定）
% 来源：樊 例 9.96，印刷页 632（PDF p641）
% ---------------------------------------------------------------------------
\begin{fdeep}{Gram 矩阵：半正定恒成立，正定 $\iff$ 线性无关}
设 $\gamma_{1},\gamma_{2},\dots,\gamma_{t}$ 是 $n$ 维欧氏空间 $V$ 中的向量组，则它\textbf{线性无关}的充分必要条件是它的 Gram 矩阵 $G(\gamma_{1},\gamma_{2},\dots,\gamma_{t})$ \textbf{正定}；更一般地，欧氏空间中\textbf{任一}向量组的 Gram 矩阵都是\textbf{半正定}矩阵。
\textbf{一句话证明}（把二次型认成向量的长度平方）：记 $A=G(\gamma_{1},\dots,\gamma_{t})$，任取 $x\in\mathbb R^{t}$ 并令 $\xi=(\gamma_{1},\gamma_{2},\dots,\gamma_{t})x$，则
\fml{x^{\mathrm T}Ax=x^{\mathrm T}G(\gamma_{1},\dots,\gamma_{t})x=(\xi,\xi)\ge0,}
即 $A$ 半正定；当 $\gamma_{1},\dots,\gamma_{t}$ 线性无关时 $x\ne0\Rightarrow\xi\ne0\Rightarrow(\xi,\xi)>0$，故 $A$ 正定。反过来，若 $A$ 正定而 $\xi=0$，则 $x^{\mathrm T}Ax=0$，由正定性得 $x=0$，即向量组线性无关。$A$ 是实对称矩阵这一点不必单独证：第 $(i,j)$ 元与第 $(j,i)$ 元都是 $(\gamma_{i},\gamma_{j})$。

【反哺点】服务考纲〔二次型〕"掌握正定二次型与正定矩阵的判别"及〔向量〕内积。这条把两章打通：\textbf{要证一个实对称矩阵正定，最省力的办法是把它写成 Gram 矩阵}（找 $\gamma_{i}$ 使 $a_{ij}=(\gamma_{i},\gamma_{j})$），然后一句 $x^{\mathrm T}Ax=|\xi|^{2}\ge0$ 收尾；反之判向量组线性无关也可转为判 $|G|>0$。
\end{fdeep}


% ---------------------------------------------------------------------------
% 【深化】服务考点：Gram 行列式（考纲〔行列式〕〔向量〕内积）
% 来源：樊 例 9.92、例 9.93，印刷页 630–631（PDF p639–p640）
% ---------------------------------------------------------------------------
\begin{fdeep}{Gram 行列式：正交化不变性与 Gram 不等式}
\textbf{结论一（施密特正交化不改变 Gram 行列式）}：设线性无关向量组 $\alpha_{1},\dots,\alpha_{m}$ 经施密特方法化成两两正交的 $\beta_{1},\dots,\beta_{m}$，则
\fml{|G(\alpha_{1},\dots,\alpha_{m})|=|G(\beta_{1},\dots,\beta_{m})|=|\beta_{1}|^{2}|\beta_{2}|^{2}\cdots|\beta_{m}|^{2}.}
证明思路：正交化过程可写成 $(\beta_{1},\dots,\beta_{m})=(\alpha_{1},\dots,\alpha_{m})Q$，其中 $Q$ 是主对角元全为 $1$ 的上三角矩阵，故 $|Q|=1$，于是 $G(\beta)=Q^{\mathrm T}G(\alpha)Q$，两边取行列式即得。
\textbf{结论二（Gram 不等式）}：对任一非零向量组，
\fml{0\le|G(\alpha_{1},\dots,\alpha_{m})|\le|\alpha_{1}|^{2}|\alpha_{2}|^{2}\cdots|\alpha_{m}|^{2},}
其中左边等号成立当且仅当 $\alpha_{1},\dots,\alpha_{m}$ 线性相关；右边等号成立当且仅当它们两两正交（两个等号不能同时成立）。

【反哺点】服务考纲〔行列式〕〔向量〕内积。记忆单元是"$|G|=$ 以这些向量为棱的平行体的体积平方"：线性相关时体积为零，两两正交时体积最大（等于各棱长之积）。有了它，$\operatorname{rank}G=\operatorname{rank}(\alpha_{1},\dots,\alpha_{m})$、线性无关判据、下一条 Hadamard 不等式就是同一件事的三种说法。
\end{fdeep}


% ---------------------------------------------------------------------------
% 【深化】服务考点：Hadamard 不等式（考纲〔行列式〕）
% 来源：樊 例 9.94，印刷页 631（PDF p640）
% ---------------------------------------------------------------------------
\begin{fdeep}{Hadamard 不等式：$|\det A|\le$ 各列长之积}
设 $A=(a_{ij})\in M_{n}(\mathbb R)$，则
\fml{|\det A|\le\prod_{j=1}^{n}\sqrt{\sum_{i=1}^{n}a_{ij}^{2}},}
等号成立当且仅当 $A$ 的列向量两两正交。证明就是把 $A$ 的列向量组记为 $\alpha_{1},\dots,\alpha_{n}$，则 $G(\alpha_{1},\dots,\alpha_{n})=A^{\mathrm T}A$，于是
\fml{(\det A)^{2}=\det(A^{\mathrm T}A)=|G(\alpha_{1},\dots,\alpha_{n})|\le\prod_{j=1}^{n}|\alpha_{j}|^{2}
=\prod_{j=1}^{n}\sum_{i=1}^{n}a_{ij}^{2},}
两边开方即可（不等号用的是上一条的 Gram 不等式）。

【反哺点】服务考纲〔行列式〕"会用行列式的性质计算行列式"。考题若问"证明某行列式的绝对值不超过若干因子之积"，标准切口就是把它认成 Gram 行列式 $|A^{\mathrm T}A|$ 再用 Gram 不等式；反过来，$|\det A|$ 最大当且仅当各列正交，这也是"正交矩阵行列式为 $\pm1$"的另一说法。
\end{fdeep}


% ---------------------------------------------------------------------------
% 【深化】服务考点：正交变换保持 Gram 矩阵（考纲〔向量〕正交变换）
% 来源：樊 例 9.95，印刷页 631–632（PDF p640–p641）
% ---------------------------------------------------------------------------
\begin{fdeep}{正交变换保持 Gram 矩阵不变}
若 $\sigma$ 是欧氏空间 $\mathbb R^{n}$ 的正交变换，记 $\beta_{i}=\sigma(\alpha_{i})$，则
\fml{G(\beta_{1},\beta_{2},\dots,\beta_{m})=G(\alpha_{1},\alpha_{2},\dots,\alpha_{m}),}
因为 $(\beta_{i},\beta_{j})=(\sigma(\alpha_{i}),\sigma(\alpha_{j}))=(\alpha_{i},\alpha_{j})$。特别地，Gram 行列式的值在正交变换下不变；若 $\alpha_{1},\dots,\alpha_{n}$ 线性无关，则正交变换不改变"这组向量的体积"。

【反哺点】服务考纲〔向量〕"了解内积""正交矩阵……的性质"。这句话是"正交变换不改变任何内积信息"的精确表述：凡是只依赖内积的量（长度、夹角、距离、Gram 行列式）在正交变换下都不变。求正交变换下的 Gram 行列式时，直接写"等于原向量组的"即可，不需重新计算。
\end{fdeep}


% ---------------------------------------------------------------------------
% 【深化】服务考点：点到子空间的距离（考纲〔向量〕内积正交）
% 来源：樊 例 9.97，印刷页 632–633（PDF p641–p642）
% ---------------------------------------------------------------------------
\begin{fdeep}{点到子空间的距离：一个纯用 Gram 行列式的公式}
设 $W$ 是 $n$ 维欧氏空间 $V$ 的子空间，$\alpha\in V$，$\alpha_{1},\dots,\alpha_{m}$ 是 $W$ 的一个基，则 $\alpha$ 到 $W$ 的距离为
\fml{d(\alpha,W)=\sqrt{\frac{|G(\alpha_{1},\alpha_{2},\dots,\alpha_{m},\alpha)|}{|G(\alpha_{1},\alpha_{2},\dots,\alpha_{m})|}}.}
\textbf{证明思路}（只须对 $\mathbb R^{n}$ 做）：令 $A=(\alpha_{1},\dots,\alpha_{m})$，则上式右端等于 $\sqrt{\alpha^{\mathrm T}\alpha-\alpha^{\mathrm T}A(A^{\mathrm T}A)^{-1}A^{\mathrm T}\alpha}$。一方面用分块初等变换把 $\begin{pmatrix}A^{\mathrm T}A&A^{\mathrm T}\alpha\\ \alpha^{\mathrm T}A&\alpha^{\mathrm T}\alpha\end{pmatrix}$ 化为 $\begin{pmatrix}A^{\mathrm T}A&0\\ 0&\alpha^{\mathrm T}\alpha-\alpha^{\mathrm T}A(A^{\mathrm T}A)^{-1}A^{\mathrm T}\alpha\end{pmatrix}$ 并取行列式；另一方面由勾股定理得 $d(\alpha,W)^{2}=|\alpha|^{2}-|\alpha'|^{2}$（$\alpha'$ 为 $\alpha$ 在 $W$ 上的正交投影），两式比较即得。

【反哺点】服务考纲〔向量〕"了解内积""正交"。这是一个\textbf{不动几何、不求投影}的距离公式：只须写出两个 Gram 行列式。内核仍是"$|G|=$ 体积平方"——分母是底面体积平方，分子是添上 $\alpha$ 后的体积平方。
\end{fdeep}


% ---------------------------------------------------------------------------
% 【深化】服务考点：正交变换存在的充要条件（考纲〔向量〕正交变换）
% 来源：樊 例 9.98，印刷页 633–634（PDF p642–p643）
% ---------------------------------------------------------------------------
\begin{fdeep}{何时存在正交变换把一组向量映到另一组向量}
设 $\alpha_{1},\alpha_{2},\dots,\alpha_{m}$ 与 $\beta_{1},\beta_{2},\dots,\beta_{m}$ 是 $n$ 维欧氏空间 $V$ 中的两组向量，则存在 $V$ 的正交变换 $\tau$ 使 $\tau(\alpha_{i})=\beta_{i}\ (i=1,2,\dots,m)$，当且仅当
\fml{(\alpha_{i},\alpha_{j})=(\beta_{i},\beta_{j}),\qquad i,j=1,2,\dots,m,}
即两组向量的 \textbf{Gram 矩阵完全相同}；注意条件\textbf{不要求}两组向量线性无关。
\textbf{证明思路}：必要性显然（正交变换保内积）。充分性分两种情形：若 $\alpha_{1},\dots,\alpha_{m}$ 线性无关，则由 $G(\alpha)=G(\beta)$ 非奇异知 $\beta_{1},\dots,\beta_{m}$ 也线性无关且 $\dim W_{1}=\dim W_{2}=m$；把 $V$ 作正交直和分解 $V=W_{1}\oplus W_{1}^{\perp}=W_{2}\oplus W_{2}^{\perp}$，取两个正交补的标准正交基对应着映过去，这样定义的 $\tau$ 保内积从而是正交变换。若向量组线性相关，则对 $m$ 用归纳法：$m=1$ 时取 $\eta=\dfrac{\alpha_{1}-\beta_{1}}{|\alpha_{1}-\beta_{1}|}$，令 $\tau(\xi)=\xi-2(\xi,\eta)\eta$ 即可；归纳步把 $\alpha_{m}$ 用前面向量线性表示再验证。

【反哺点】服务考纲〔向量〕"了解内积""正交矩阵……的性质"。它是 §9.2 与 §9.4 的交汇点：\textbf{内积（Gram 矩阵）是正交变换下的完全不变量}——两组向量能被同一个正交变换互相搬运，当且仅当内积表相同。凡是"是否存在正交变换使……"的存在性问题，写出并比较 Gram 矩阵即可。
\end{fdeep}


% ---------------------------------------------------------------------------
% 以下按深化提取原则跳过
% ---------------------------------------------------------------------------
% fan p604 §9.2.2 第（5）条：$\mathscr{A}$ 在标准正交基下的矩阵为
%   $\operatorname{diag}(E_r,-E_s,A_1,\dots,A_t)$（$A_i$ 为二阶旋转块）——
%   这是正交变换的完整分类理论（标准形），超出考研大纲，按 prompt 跳过。
%   （上面"正交 + 特征值全实 $\Rightarrow$ 对称"一块只在证明中引用了它的形式，
%     未展开分类本身的推导。）
% fan p604–605 例 9.30（给定两个向量求 4 阶正交阵使其为前两列）：纯构造题，
%   可复用方法已并入「求正交变换的三条现成路径」（路径一）。
% fan p606 例 9.33（南京大学 2015 等，非奇异矩阵的 $PAQ$ 与奇异值分解）：超纲，跳过。
% fan p606–607 例 9.34：结论已提取，其"仅在 $\mathbb R^n$ 中讨论"的对照证明不另立块。
% fan p608 例 9.37（三阶正交阵写成旋转与镜面反射之积 $A=CR$）：正交变换分类，超纲，跳过。
% fan p610–612 例 9.42、例 9.43（可逆矩阵的正交三角分解 $A=QR$ 及唯一性、
%   Cholesky 分解、具体 4 阶数值分解）：QR 与 Cholesky 属超纲深化内容，跳过。
%   （例 9.42 唯一性证明用到的那条引理"上三角正交阵必为对角阵且对角元为 $\pm1$"
%     已在上面单独成块提取。）
% fan p612 例 9.45 的完整解答：结论（1）（2）（3）已提取，取 $B=E_{ii}$ 验算的细节不另立块。
% fan p613–614 例 9.46（任一正交阵可表为 Givens 旋转矩阵之积）及其逐列消元构造：
%   Householder/Givens 深化内容，超纲，跳过。
% fan p615–616 例 9.49、例 9.50：解答中的三段构造已提炼为
%   「求正交变换的三条现成路径」，具体数值验算不进正文。
% fan p637 上半页：§9.3 末尾一个内积等式的证明，属 §9.3 且为证明细节，跳过。
% fan p639 例 9.91（北京大学 2000，(0,1) 矩阵行和 $r$、行内积 $m$ 时
%   $|AA^{\mathrm T}|=[r+(s-1)m](r-m)^{s-1}$、$s\le n$、$AA^{\mathrm T}$ 正定）：
%   其内核"$AA^{\mathrm T}$ 是行向量组的 Gram 矩阵"与正定性已由 Gram 两块覆盖，
%   针对 (0,1) 矩阵的具体算式不再单列。
% ===========================================================================


% ===========================================================================
%  【主控裁决 · 第 6 章 deep_ch06_fan 取舍】（依据 AGENTS.md §2 决定二）
%
%  考纲对〔向量〕这一块的要求层次是「**了解**规范正交基、正交矩阵的概念以及它们的性质」，
%  故：正交矩阵的**基本性质与判定**收，正交变换的**深层理论**不收。
%
%  【保留 8 块】
%   ✓ 块5　{$|Q|=\pm1$ 的两条直接推论
%   ✓ 块8　{正交矩阵的元素与代数余子式：$a_{ij}=\pm A_{ij}$
%   ✓ 块9　{上（下）三角的正交矩阵必为对角矩阵，对角元只能是 $\pm1$
%   ✓ 块12　{判定补充二：元素平方和 $\sigma(A)=\operatorname{tr}(A^{\mathrm T}A
%   ✓ 块14　{Gram 矩阵（度量矩阵）：定义与三种来源
%   ✓ 块15　{Gram 矩阵的秩 $=$ 向量组的秩（线性无关判据）
%   ✓ 块16　{Gram 矩阵：半正定恒成立，正定 $\iff$ 线性无关
%   ✓ 块17　{Gram 行列式：正交化不变性与 Gram 不等式
%
%  【删除及理由】
%   ✗ 块1　{正交变换：四条等价刻画（把矩阵语言升级为变换语言）
%   ✗ 块2　{保持内积的变换自动是线性的
%   ✗ 块3　{保距变换还差一个"保原点"
%   ✗ 块4　{正交矩阵的特征值模恒为 $1$
%   ✗ 块6　{三阶正交矩阵的特征多项式必为 $\lambda^{3}-a\lambda^{2}+a\lambda-1
%   ✗ 块7　{三阶正交矩阵的迹：一个恒等式与旋转角公式
%   ✗ 块10　{正交矩阵的特征值全为实数 $\iff$ 它是对称矩阵
%   ✗ 块11　{判定补充一：$(A+I)^{-1}+(A^{\mathrm T}+I)^{-1}=I$
%   ✗ 块13　{求正交变换的三条现成路径
%   ✗ 块18　{Hadamard 不等式：$|\det A|\le$ 各列长之积
%   ✗ 块19　{正交变换保持 Gram 矩阵不变
%   ✗ 块20　{点到子空间的距离：一个纯用 Gram 行列式的公式
%   ✗ 块21　{何时存在正交变换把一组向量映到另一组向量
%
%  归类理由：
%   · 正交变换深层理论（超纲，考纲只要"了解概念及性质"）：
%       块1 正交变换四条等价刻画、块2 保内积者自动线性、块3 保距还差保原点、
%       块6 三阶正交阵特征多项式、块7 迹与旋转角公式、块13 求正交变换三条路径、
%       块21 何时存在正交变换把一组向量映到另一组
%   · 竞赛级/过专（非考纲主线）：块10 特征值全实 ⟺ 对称、块11 $(A+I)^{-1}$ 恒等式、
%       块18 Hadamard 不等式、块20 点到子空间距离的 Gram 行列式公式
%   · 与本章他片重复：块19 正交变换保持 Gram 矩阵不变（与块14-17 的 Gram 内容重叠）
%   · 已被主控手写块覆盖：正交矩阵六条等价刻画、保长度（deep_ch06_authored）
%   · 已被 deep_ch06 覆盖：度量矩阵的换基合同（块14 的定义部分与其协调，保留其"三种来源"）
%
%  保留 8 块 = 正交矩阵基本性质与判定（4 块）+ Gram 矩阵完整理论（4 块）。
% ===========================================================================
